% !TEX root = ../main.tex
\section{Grounding, Context, and the Taxonomy of Assertion}
\label{sec:grounding}

\subsection{The grounding operator}
\label{sec:grounding-operator}

The negative result of Sections~\ref{sec:output-operator}
and~\ref{sec:hyperintensionality} might suggest that no epistemically
interesting structure survives in machine output. It does not follow, and the
rest of this paper develops the structure that remains. The move is to stop
asking whether $M$'s assertions express knowledge and start asking whether they
are \emph{grounded}.

\begin{definition}[Grounding operator]
\label{def:grounding}
$\Grd_M(\phi \mid c)$ holds iff $\phi$ is supported by $M$'s evidential base $E$
relative to practical context $c$---that is, iff
\begin{enumerate}
  \item $E$ contains information that evidentially favours $\phi$, and that
        information is \emph{applicable} at $c$: authentic, current at $\tau$,
        and of the right jurisdiction $j$ and intended use $u$; and
  \item the mechanism producing the assertion of $\phi$ is counterfactually
        sensitive to that information---roughly, had $E$ not supported $\phi$,
        the mechanism would not, so far as this content is concerned, have
        produced it.
\end{enumerate}
\end{definition}

Here $E$ comprises the training corpus, any retrieval-provided context, and the
parametric encodings derived from them. Four clarifications, because everything
in Section~\ref{sec:hallucination} turns on this operator.

\paragraph{Support is epistemic, not causal.} Every output of $M$ is trivially
caused by training, and that cannot be what $\Grd_M$ tracks; if it were, the
operator would be satisfied by everything and the account would be empty.
Condition (2) is what does the work. It excludes the case where $E$ happens to
support $\phi$ but the output was generated by formulation-level statistics
indifferent to that support. Attribution and interpretability methods are, in
effect, attempts to operationalise condition (2), a point
Section~\ref{sec:evaluation} develops along with the limits of what they can
currently deliver.

\paragraph{Grounding is not factive.} $E$ may contain errors; a corpus saturated
with a widespread misconception supports the misconception. Grounding is a
relation between an assertion and a body of evidence, and bodies of evidence can
be bad. This non-factivity is not a defect of the definition. It is the
definition's point, and Section~\ref{sec:error} turns on it.

\paragraph{The context index is load-bearing.} Condition (1) requires
applicability at $c$, and this is the addition that
Section~\ref{sec:contextual-truth} was setting up. Without it, a claim supported
anywhere in $E$ counts as grounded, and the operator loses its grip on the
failure modes that dominate regulated practice. A regulatory position that was
authoritative in 2019 and superseded in 2024 supports an assertion made in 2019
and does not support the same assertion made now. Evidence drawn from one
jurisdiction's regime does not support a claim advanced for a decision under
another's. The unindexed operator cannot say this; the indexed one can, and
Section~\ref{sec:misapplied} shows that saying it identifies a failure class with
its own remedy.

\paragraph{Both operators are hyperintensional and admit of degree.} $\Grd_M$ is
no more immune to formulation-sensitivity than $\Out_M$ is, and support comes in
degrees. I work with a thresholded, formulation-relative version throughout for
uniformity, and nothing below trades on the idealisation.

The pair $(\Out_M, \Grd_M)$ replaces the pair (belief, justification) for systems
of this kind. $\Out_M$ is what the system does; $\Grd_M$ is whether the doing is
evidentially connected to anything applicable.

\subsection{The taxonomy}
\label{sec:taxonomy}

Crossing grounding with truth yields four exhaustive and exclusive categories of
asserted content:

\begin{table}[ht]
  \centering
  \begin{tabular}{@{}lll@{}}
    \toprule
     & $\Grd_M(\phi \mid c)$ & $\neg\,\Grd_M(\phi \mid c)$ \\
    \midrule
    $\phi$ true
      & (i) veridical transmission
      & (ii) veridical hallucination \\
    $\phi$ false
      & (iii) error
      & (iv) fabrication \\
    \bottomrule
  \end{tabular}
  \caption{Grounded against true. Cells (i) and (iv) are the diagonal ordinary
    talk handles well: getting it right for the right reason, and making it up
    and getting it wrong. The off-diagonal cells are where the truth-conditional
    conception of hallucination collapses.}
  \label{tab:taxonomy}
\end{table}

Because grounding is now indexed, the grounded column admits a further
distinction that the unindexed taxonomy cannot draw. An assertion may be
supported by evidence in $E$ that is inapplicable at $c$---stale, wrong
jurisdiction, wrong intended use. Such an assertion is ungrounded relative to $c$
by Definition~\ref{def:grounding}, and so falls in the right-hand column
alongside fabrication, despite the presence of supporting text in $E$. That
placement is deliberate and I defend it in Section~\ref{sec:misapplied}: what
matters for a consequential decision is whether the assertion is supported by
evidence that bears on \emph{this} decision, and evidence that does not bear on
it provides no support at all, however authentic it was in its own context.
